\chapter{A játékosháló felépítése, keresése és mérései}

\section{A gráf felépítése}

Az elsődleges gráfállapotot a Rust-réteg határozza meg. A Python oldal megőrzi az adatgyűjtés, a gyors prototipizálás és a kompatibilitási gráfgenerálás szerepét. A \textit{flexset} és \textit{soloq} összehasonlítása, az asszortativitás és a gráfvizualizáció Rust-alapú gráfmodellre támaszkodik.

\subsection{Kapcsolataggregáció}

A Rust gráfmodell minden játékospárra külön őrzi a szövetségesi és ellenfél-oldali bizonyítékot:

\begin{equation}
w_{ally}(u,v)=\sum_{m \in M}\mathbf{1}[\exists i \in \{1,2\}: u \in T_{m,i} \land v \in T_{m,i}]
\end{equation}
\begin{equation}
w_{enemy}(u,v)=\sum_{m \in M}\mathbf{1}[\exists i \neq j: u \in T_{m,i} \land v \in T_{m,j}]
\end{equation}

A domináns reláció \textit{ally}, ha \textit{allyWeight} $\geq$ \textit{enemyWeight}, különben \textit{enemy}. A gráf eredményei elsősorban asszociatív játékosgráfként, nem szó szerinti barátsági vagy ellenségességi hálóként értelmezendők.

\subsection{A Rust gráfállapot}

A Rust-réteg a gráfot közvetlenül a meccs-JSON-állományokból és a játékos-adatbázisból építi fel. A \textit{GraphState} a szomszédsági listák mellett tárolja a \textit{pair\_relations} szerkezetet, a játékosprofilok gyors elérését, a klasztertagságot, a landmark-távolságokat, a minimumköltségeket és az aktív datasethez tartozó snapshot-metaadatot.

\subsection{Szűrés és a kutatási gráf}

A kutatási gráf láthatósági küszöbe \textit{totalMatches} $\geq 2$: egyetlen közös mérkőzés még nem elég ahhoz, hogy egy kapcsolat stabil élként jelenjen meg. A klaszterezés fő iránya a Rust-oldali, determinisztikus \textit{bounded-ally-groups} logika. A vizuális klaszterek mérete szándékosan korlátozott, legfeljebb körülbelül húsz játékosra.

\subsection{A gráf artifact-rétege}

A Rust gráf-export pipeline előkészített, datasethez kötött cache-artifaktumokat ír ki: a \textit{manifest.json} rögzíti a datasetet és a gráfépítő verziót; a bináris fájlok (\textit{node\_positions.f32}, \textit{edge\_pairs.u32}) a Three.js/WebGL réteg bemeneti állományai; a \textit{cluster\_meta.json} a klaszterazonosítókat tartalmazza. Ez az artifact-réteg tekintendő a kutatási vizualizáció fő igazságforrásának.

\section{Útvonalkeresés és vizualizáció}

\subsection{Keresési algoritmusok}

A jelenlegi Rust futtatókörnyezet négy algoritmust támogat: BFS, Dijkstra, kétirányú keresés és A*. Az algoritmusok két értelmezésben futhatnak: a \textit{social-path} csak szövetségesi kapcsolatokon, a \textit{battle-path} pedig az összes ismétlődő kapcsolaton.

Az élköltség-modell: súlyozatlan módban minden él költsége fix $1000$. Súlyozott módban:
\begin{equation}
c(u,v)=\left\lceil \frac{1000}{w(u,v)} \right\rceil
\end{equation}
ahol $w(u,v)$ az adott útvonalmódnak megfelelő kapcsolaterősség.

Az \ref{tab:pathfinder-algorithms} táblázat a Pathfinder Labban a négy keresési algoritmus jellemzőit hasonlítja össze.

\begin{table}[H]
\centering
\caption{A Pathfinder Labban szereplő algoritmusok összevetése}
\label{tab:pathfinder-algorithms}
\begin{tabularx}{\textwidth}{|>{\raggedright\arraybackslash}p{2.8cm}|>{\raggedright\arraybackslash}p{2.8cm}|X|}
\hline
\textbf{Algoritmus} & \textbf{Mikor természetes?} & \textbf{Megjegyzés a projektben} \\
\hline
BFS & súlyozatlan legrövidebb út & Jó baseline, könnyen magyarázható \\
Dijkstra & súlyozott utak & Referenciaalgoritmus a költség-alapú módokban \\
Kétirányú keresés & nagyobb gráf, kevésbé irányított keresés & Gyakorlati gyorsítás bizonyos esetekben \\
A* & jól kialakított alsó becsléssel & A projektben egzakt marad, nem csak vizuális heurisztika \\
\hline
\end{tabularx}
\end{table}

\subsection{Egzakt A* és ALT heurisztika}

Az ALT (A* + Landmarks + Triangle inequality) heurisztika \cite{goldberg2005} elvéből indul ki: ha $L$ egy landmark-csomópont, akkor
\begin{equation}
d(u, t) \geq |d(L, t) - d(L, u)|
\end{equation}
minden $L$ landmarkra. A landmark-távolságok a \textit{GraphState} felépítésekor előre kiszámítódnak minden \textit{ModeKey}-kombinációra, és az összes keresési kérésre újrafelhasználhatók. Az A* a helyes legrövidebb utat adja; egységtesztek igazolják, hogy az A* eredménye minden releváns esetben egyezik a Dijkstra-eredménnyel.

\subsection{Pathfinder Lab}

A Pathfinder Lab a négy keresési algoritmus interaktív tesztkörnyezete: az eredmények lépésenként visszajátszhatók, a két útvonalmód (social-path, battle-path) közvetlenül összehasonlítható.

A \ref{fig:pathfinder-replay} ábra a Pathfinder Lab replay panelét mutatja, amelyen az útvonalösszegző, a léptetési vezérlők és a bejárt csomópontok statisztikái láthatók.

\begin{figure}[H]
\centering
\includegraphics[width=0.98\textwidth,height=0.74\textheight,keepaspectratio]{pathfinder-replay-panel.png}
\caption{A Pathfinder Lab replay panele: az útvonalösszegző, a léptetési vezérlők és a meglátogatott csomópontok statisztikái}
\label{fig:pathfinder-replay}
\end{figure}

\begin{table}[H]
\centering
\caption{BFS és súlyozott keresés összehasonlítása: ugyanazon játékospár, eltérő közelségi fogalmak}
\label{tab:bfs-weighted-compare}
\begin{tabularx}{\textwidth}{|>{\raggedright\arraybackslash}p{3cm}|>{\raggedright\arraybackslash}p{3cm}|>{\raggedright\arraybackslash}p{2cm}|X|}
\hline
\textbf{Algoritmus} & \textbf{Metrika} & \textbf{Eredmény} & \textbf{Értelmezés} \\
\hline
BFS & Minimális ugrásszám & 3 & Legrövidebb szociális lánc \\
Dijkstra & Min. súlyozott költség & 8 & Legerősebb szövetségesi útvonal \\
A* & Min. súlyozott költség & 8 & Egyezik Dijkstrával \\
Kétirányú & Min. ugrásszám & 3 & Egyezik BFS-sel \\
\hline
\end{tabularx}
\end{table}

A BFS által talált 3-ugrásos útvonal strukturálisan létezik, de gyenge vagy egyszeri meccsből eredő kapcsolatokon halad át. A Dijkstra súlyozott módban elkerüli ezeket, és helyette egy 8-ugrásos útvonalat talál, amelyen minden él ismétlődő szövetségesi viszony.

Az \ref{tab:algo-comparison} táblázat a négy algoritmus teljesítményjellemzőit és fő erősségeit hasonlítja össze.

\begin{table}[H]
\centering
\caption{A négy algoritmus összehasonlítása a legfontosabb dimenziók mentén}
\label{tab:algo-comparison}
\begin{tabularx}{\textwidth}{|>{\raggedright\arraybackslash}p{3cm}|>{\raggedright\arraybackslash}p{2.5cm}|>{\raggedright\arraybackslash}X|>{\raggedright\arraybackslash}X|}
\hline
\textbf{Algoritmus} & \textbf{Optimális?} & \textbf{Memória} & \textbf{Fő erősség a projektben} \\
\hline
BFS & Igen (súlyozatlan) & $O(n)$ & Egyszerű baseline, jól magyarázható replay \\
Dijkstra & Igen (súlyozott) & $O(n + m)$ & Referencia-algoritmus, lexikografikus döntetlen-kezelés \\
Kétirányú & Igen (súlyozatlan) & $O(n)$ & Potenciális gyorsulás klaszteres gráfon \\
A* & Igen (heurisztika elfogadható) & $O(n)$ & Célirányos bejárás, egzakt eredmény \\
\hline
\end{tabularx}
\end{table}

\subsection{A 3D globális vizualizáció}

A \textit{bird's-eye sphere} a teljes hálózat vizuális áttekintésére szolgál. A Rust-réteg előre exportálja az összes fontos artefaktumot, a frontend pedig csak renderel.

Az \ref{tab:birdseye-artifacts} táblázat a bird's-eye export artifaktumait és azok szerepét mutatja be.

\begin{table}[H]
\centering
\caption{A bird's-eye export fő artifaktumai}
\label{tab:birdseye-artifacts}
\begin{tabularx}{\textwidth}{|>{\raggedright\arraybackslash}p{4.2cm}|X|}
\hline
\textbf{Fájl} & \textbf{Szerep} \\
\hline
\textit{manifest.json} & Metaadatok, darabszámok, skálázási és generálási információk \\
\textit{node\_meta.json} & Csomópontszintű leíró metaadatok, frontend-hozzárendelések \\
\textit{node\_positions.f32} & Lebegőpontos koordináták tömbje a gyors rendereléshez \\
\textit{node\_metrics.u32} & Tömörített numerikus csomóponti mutatók \\
\textit{edge\_pairs.u32} & Él-végpontpárok tömbje \\
\textit{edge\_props.u32} & Élhez tartozó kódolt tulajdonságok \\
\hline
\end{tabularx}
\end{table}

A \ref{fig:graph-artifact-dataflow} ábra a gráf artifact-csomag teljes adatfolyamát mutatja be a Rust-exportőrtől a Three.js renderelőig.

\begin{figure}[H]
\centering
\makebox[\textwidth][c]{\includegraphics[width=1.3\textwidth]{graph artifact package dataflow.pdf}}
\caption{A gráf artifaktumcsomag adatfolyama: a Rust-exportőr bináris és JSON-rétegeket ír, az Express API kiszolgálja, a frontend WebGL-renderelőbe fogyasztja}
\label{fig:graph-artifact-dataflow}
\end{figure}

A \ref{fig:graph-artifact-components} ábra az artifact-csomag bináris és JSON-fájljainak szerkezetét mutatja be.

\begin{figure}[H]
\centering
\hspace*{1cm}\includegraphics[width=1.8\textwidth]{graph artifact components.pdf}
\caption{A gráf-export pipeline bináris és JSON artifact-fájljainak szerkezete}
\label{fig:graph-artifact-components}
\end{figure}

A \ref{fig:birds-eye} ábra a teljes \textit{flexset} hálózat 3D madártávlati nézetét mutatja, az előre generált Rust-artifaktumok alapján renderelve.

\begin{figure}[H]
\centering
\includegraphics[width=0.98\textwidth,height=0.72\textheight,keepaspectratio]{birds-eye-current-state.png}
\caption{A teljes hálózat 3D madártávlati nézete az előre generált Rust-artifaktumokra építve}
\label{fig:birds-eye}
\end{figure}

A \ref{fig:flexset-clusters} ábrán a \textit{flexset} gráf sűrűbb klaszterfoltjai, valamint a térbeli elkülönítés látható a globális nézetben.

\begin{figure}[H]
\centering
\includegraphics[width=0.98\textwidth,height=0.72\textheight,keepaspectratio]{flexset_associative_graph_clusters.png}
\caption{Sűrűbb klaszterfoltok és a térbeli elkülönítés a globális nézetben}
\label{fig:flexset-clusters}
\end{figure}

\section{Asszortativitás}

\subsection{A számítási definíció}

Az implementáció numerikus asszortativitást számol Pearson-korreláció formájában:

\begin{equation}
r = \frac{\sum_i (x_i-\bar{x})(y_i-\bar{y})}{\sqrt{\sum_i (x_i-\bar{x})^2}\sqrt{\sum_i (y_i-\bar{y})^2}}
\end{equation}

Az irányítatlan éleket az implementáció mindkét endpoint-orientációval figyelembe veszi. A gráfmód alapján az él-súlyozás:
\begin{equation}
w(u,v) = \begin{cases}
\textit{ally\_weight} & \text{ha } \textit{social-path} \\
\textit{total\_matches} & \text{ha } \textit{battle-path}
\end{cases}
\end{equation}

\subsection{Asszortativitási eredmények: flexset}

Az \ref{tab:assort-flexset} táblázat a \textit{flexset} asszortativitási eredményeit tartalmazza kétféle gráfmódban.

\begin{table}[H]
\centering
\caption{Asszortativitás: \textit{flexset}}
\label{tab:assort-flexset}
\begin{tabular}{|l|l|r|r|}
\hline
\textbf{Gráfmód} & \textbf{Metrika} & \textbf{Korreláció} & \textbf{Mintaszám} \\
\hline
\textit{social-path} & \textit{opscore} & 0.2346 & 61\,938 \\
\textit{social-path} & \textit{feedscore} & 0.1816 & 61\,938 \\
\textit{battle-path} & \textit{opscore} & 0.1461 & 179\,492 \\
\textit{battle-path} & \textit{feedscore} & $-0.0108$ & 179\,492 \\
\hline
\end{tabular}
\end{table}

A \ref{fig:flexset-assortavity} ábra a \textit{flexset} social-path asszortativitási elemzésének vizualizációját mutatja \textit{opscore} és \textit{feedscore} szerint.

\begin{figure}[H]
\centering
\includegraphics[width=0.98\textwidth,height=0.74\textheight,keepaspectratio]{f-assortavity-social-path-opfeedscore.png}
\caption{\textit{flexset} social-path asszortativitás \textit{opscore} és \textit{feedscore} szerint. 15\,421 jogosult csomópont és 61\,938 elemzett él; globális együttható \textit{opscore} esetén 0.235, \textit{feedscore} esetén 0.182.}
\label{fig:flexset-assortavity}
\end{figure}

\subsection{Asszortativitási eredmények: soloq}

Az \ref{tab:assort-soloq} táblázat a \textit{soloq} asszortativitási eredményeit tartalmazza, összehasonlítva a két gráfmódot.

\begin{table}[H]
\centering
\caption{Asszortativitás: \textit{soloq}}
\label{tab:assort-soloq}
\begin{tabular}{|l|l|r|r|}
\hline
\textbf{Gráfmód} & \textbf{Metrika} & \textbf{Korreláció} & \textbf{Mintaszám} \\
\hline
\textit{social-path} & \textit{opscore} & 0.1750 & 37\,976 \\
\textit{social-path} & \textit{feedscore} & 0.1012 & 37\,976 \\
\textit{battle-path} & \textit{opscore} & 0.1625 & 84\,393 \\
\textit{battle-path} & \textit{feedscore} & $-0.0366$ & 84\,393 \\
\hline
\end{tabular}
\end{table}

A \ref{fig:soloq-assortavity} ábra a \textit{soloq} social-path asszortativitási eredményeit mutatja.

\begin{figure}[H]
\centering
\includegraphics[width=0.98\textwidth,height=0.74\textheight,keepaspectratio]{sq-assortavity-social-path-opfeedscore.png}
\caption{\textit{soloq} social-path asszortativitás \textit{opscore} és \textit{feedscore} szerint. 6\,768 jogosult csomópont és 37\,976 elemzett él; globális együttható \textit{opscore} esetén 0.175, \textit{feedscore} esetén 0.101.}
\label{fig:soloq-assortavity}
\end{figure}

\subsection{Értelmezés}

Mindkét datasetben pozitív asszortativitás jelenik meg az \textit{opscore} mentén a \textit{social-path} módban, ami azt jelzi, hogy a szövetségesi kapcsolatokra épülő gráfban az összekapcsolt játékosok teljesítményprofilja nem véletlenszerűen keveredik. A \textit{feedscore} viselkedése jóval érzékenyebb a gráfmódra: social-path módban mindkét dataseten mérsékelt pozitív jel látható, de battle-path módban közel nullára esik.

Az \textit{opscore} minden vizsgált módban pozitív marad, a \textit{feedscore} nem; ez a kompozit mutató erőssége az egyszerűbb halálközpontú képlettel szemben.

\section{Betweenness centralitás}

\subsection{A Brandes-algoritmus}

A Brandes-algoritmus \cite{brandes2001} a hagyományos $O(n^3)$ komplexitású betweenness-számítást $O(nm)$-re csökkenti. Forráscsomópontonként: Dijkstra/BFS futtatás, majd visszaterjesztési lépésben:

\begin{equation}
\delta[s][v] \mathrel{+}= \frac{\sigma[s][v]}{\sigma[s][w]} \cdot (1 + \delta[s][w])
\end{equation}

Súlyozott módban az élköltség:
\begin{equation}
c(u,v) = \left\lceil \frac{10^6}{\max(1, w(u,v))} \right\rceil
\end{equation}

\subsection{Párhuzamos végrehajtás}

A Rust \textit{rayon} könyvtára automatikus munkalopásos párhuzamosságot tesz lehetővé. A centralitás-modul a forráscsomópontok halmazát $k = \min(n, 4 \cdot \text{szálak száma})$ egyenlő méretű csomóponthalmazra osztja. A forráscsomópontok rendezése és a chunk-határok rögzítése garantálja, hogy azonos input esetén mindig ugyanaz a kimenet.

\subsection{Eredmény}

A \ref{fig:brandes-result} ábra a Brandes-algoritmus futáseredményének képernyőjét mutatja, beleértve a párhuzamos és soros futásidők összehasonlítását.

\begin{figure}[H]
\centering
\includegraphics[width=0.98\textwidth,height=0.68\textheight,keepaspectratio]{brandes-result.png}
\caption{Brandes betweenness centralitás eredményképernyő. 15\,421 runtime csomópontból 1\,716 projektált csomópont és 2\,766 elemzett él. Párhuzamos futás: 6.38 ms; soros baseline: 60.85 ms; gyorsulás: 9.54-szeres; 16 Rayon szál; soros-párhuzamos delta: 0.000000.}
\label{fig:brandes-result}
\end{figure}

A \ref{fig:brandes-brokers} ábra a centralitás szempontjából kiemelkedő, közvetítő szerepű játékosokat mutatja a gráfban, sárgával jelölve.

\begin{figure}[H]
\centering
\includegraphics[width=\textwidth,height=0.78\textheight,keepaspectratio]{graph-brandes-brokerinyellow.png}
\caption{A betweenness centralitás alapján azonosított broker játékosok a gráfban, sárgával kiemelve.}
\label{fig:brandes-brokers}
\end{figure}
